\section{Related work}\label{sec:related}

We review four bodies of work this study sits between. The review is based on the sources listed in the reference list, compiled and checked between July and August 2026; it is a directed review, not a systematic one, and the statements of novelty below are made in that light.

\subsection{Digitisation of paper and photographed ECGs}
Early digitisation methods used classical image processing to remove the grid, threshold the trace and rescale pixels to millivolts \citep{ref6,ref7,ref8,ref9}. Recent work uses deep networks for trace segmentation and layout recognition \citep{ref10,ref11}. The open-source Open-ECG-Digitizer \citep{ref5}, used here, combines a segmentation network with a lead-name network and a layout-matching step, and reports validation on more than 37{,}000 images. ECGtizer \citep{ref3} and the entry that won the 2024 George B.\ Moody PhysioNet Challenge \citep{ref53} are independent alternatives; the Challenge established signal-to-noise ratio against ground truth as its standard accuracy measure \citep{ref2}. ECG-Image-Kit \citep{ref1} generates realistic printed ECG pages, optionally with artefacts, from digital records, and ECG-image-database \citep{ref4} provides paired images and signals with real-world artefacts such as perspective skew, wrinkles, stains and uneven lighting. Performance on photographs is markedly lower than on scans: Open-ECG-Digitizer reports 19.65\,dB on scans but about 12.2 and 10.5\,dB on two phone cameras \citep{ref5}, and the Challenge winner, at 12.15 on the scans-only leaderboard, scored negative signal-to-noise ratio on the mobile-phone photographs of the hidden set \citep{ref53}. Two recent studies are closest to ours: one benchmarks digitisers separately for each of nine capture or damage types of a Challenge dataset and reports signal-to-noise ratio only \citep{ref54}; the other adds a quality-assurance module after digitisation and scores it on reconstruction fidelity, not clinical error \citep{ref55}. We differ in grading each damage type over five severities, in measuring clinical values rather than signal-to-noise ratio, and in defining and counting false reassurance. We use ECG-Image-Kit for rendering and treat the digitiser as a component under test, not a contribution, and not as a solved problem for phone photographs.

\subsection{Signal and image quality assessment}
A separate literature asks whether an ECG is good enough to analyse at all. Deep classifiers separate interpretable from non-interpretable recordings with high reported accuracy \citep{ref12}, and the problem has been studied for wearable and portable devices \citep{ref13,ref15,ref16,ref19,ref20}, for false-alarm reduction \citep{ref17,ref18}, with unsupervised learning \citep{ref14,ref21} and with attention networks \citep{ref22}. Signal-quality gating before analysis of a \emph{waveform} is therefore feasible and established. The Challenge organisers likewise observe that signal-to-noise ratio does not directly reflect the clinical measurements that influence ECG interpretation \citep{ref2}. That observation motivates the present focus on error at the level of the extracted clinical value.

\subsection{Uncertainty, abstention and deferral}
When should an automated medical system decline to answer? Uncertainty-aware abstention has been studied for medical text \citep{ref23}, conformal and cost-aware deferral for clinical triage \citep{ref24,ref27}, guided deferral for language-model systems \citep{ref25}, uncertainty quantification for ECG interpretation \citep{ref30,ref31}, selective classification for prehospital risk prediction \citep{ref32}, and uncertainty estimation in other diagnostic settings \citep{ref28,ref29}. A cautionary result is that abstention is not automatically safer: when an artificial-intelligence system abstains, clinicians' error patterns shift, with false positives falling and false negatives rising \citep{ref26}. This body of work presupposes an uncertainty estimate; it does not, in the sources we reviewed, connect that estimate to a measured rate of a specific harm caused by a specific, controllable input condition.

\subsection{Patient-facing explanation and medical AI safety}
Language models are increasingly used to translate clinical output into patient-readable text, and their quality and safety have been evaluated for cardiology procedures \citep{ref33}, radiology reports \citep{ref38}, personalised plain-language summaries \citep{ref34}, chronic-disease self-management \citep{ref39} and general patient questions \citep{ref40,ref41,ref42,ref43}. A clinical-safety benchmark defines 30 metrics for medical language models \citep{ref35}. ECG-specific language-model work exists \citep{ref36,ref37} but targets clinicians interpreting a trace, not a patient reading an explanation of their own photograph.

\subsection{Gap}
Each thread stops at a boundary: digitisation stops at the signal, quality assessment at whether the signal is usable, abstention at an assumed uncertainty score, and patient-facing generation at the text. In the sources we reviewed, none measures how far a \emph{controllable image degradation} propagates into a \emph{clinically meaningful measurement} and at what severity it becomes a falsely reassuring one. That measurement is the subject of this paper. It is a necessary input to any later safeguard, whether an upload-time quality gate or an uncertainty-qualified explanation.
